\section{Nonvanishing, normalization and the final family}\label{app:endpoints}
The uses of the coefficient and address propositions have the following
order. First apply only their even-source case, whose depth is already
fixed by Proposition~\ref{ex:app-even-phase}. This gives the rough
companion and the genus lower bound. These in turn establish the odd
minimum and its finite-stage inequalities. Only then apply the odd
case of Proposition~\ref{ex:app-coefficient-rules} and the two-symbol
address construction. In particular the proof of the rough companion
does not use the later odd minimum.
We give the choice order underlying (3.3)--(3.9). This also separates
the two valuation arguments: the even upper transfer uses the common
Kato scalar, while the final contradiction uses the common Heegner
coordinate of Section~\ref{sec:limits}.

\subsection{An absolute bound of five new primes}
The theorem of Hoffstein--Luo \cite[p.~436]{hl} supplies fundamental
discriminants with nonzero central value and at most four prime
factors. Its finite splitting set is arbitrary. The bound four is
absolute; no uniform size threshold for those discriminants is needed.

\begin{proposition}[Nonvanishing with five primes and finite local prescriptions]\label{ex:app-five-prime}
Let $A/\Q$ be a fixed modular elliptic curve for this application,
let $S$ contain $2,3$ and all its bad primes, and let $B$ be a finite
prime exclusion set. In every nonempty signed unit-squareclass filter
at $S$ of even root number, there is a signed odd squarefree parameter
$u$, prime to $S\cup B$, with $L(A^u,1)\ne0$ and $\omega(u)\le5$.
Only the intermediate parameter supplied by Hoffstein--Luo is required
to be fundamental; the final field discriminant may be $4u$.
The number five is independent of the prescribed finite set. No
uniform bound on the sizes of the selected primes is asserted.
\end{proposition}
\begin{proof}\leavevmode
Let $A/\Q$ be any fixed modular elliptic curve for this application,
$S$ contain two, three and all its bad primes, and $B$ be any finite
exclusion set. Fix a nonempty signed unit filter of even root number.
By CRT and Dirichlet choose one signed prime $q_0$ in its classes,
outside $S\cup B$. Apply the theorem to the actual curve $A^{q_0}$,
putting
\[
                 S'=S\cup B\cup\{|q_0|\}\cup
                               \{p:p\mid N_{A^{q_0}}\}
\]
in its splitting set. The theorem gives a fundamental $d$ with
$\chi_d(p)=1$ for all $p\in S'$ and at most four prime factors.
The functional equation forces $d>0$, since $A^{q_0}$ has even
root number and every conductor prime splits. Splitting at two
gives $d\equiv1\pmod8$. Thus $q_0d$ has all original unit
classes, is prime to $B$, is squarefree, and has at most five
prime factors. If $q_0d\equiv3$ or $7\pmod8$, its field
discriminant is $4q_0d$; only $d$ was required to be fundamental
in the theorem. This is compatible with the coefficient
construction (E.2).
\end{proof}

